\begin{problem}[第35届全国中学生物理竞赛复赛][磁阱束缚]{91}
四、Ioffe-Pritchard 磁阱可用来束缚原子的运动，其主要部分如图所示。四根均通有恒定电流 $I$ 的长直导线 $1$、$2$、$3$、$4$ 都垂直于 $x-y$ 平面，它们与 $x-y$ 平面的交点是边长为 $2a$、中心在原点 $O$ 的正方形的顶点，导线 $1$、$2$ 所在平面与 $x$ 轴平行，各导线中电流方向已在图中标出。整个装置置于匀强磁场 $B_{0}=B_{0}k$ （$k$ 为 $z$ 轴正方向单位矢量）中。已知真空磁导率为 $\mu_{0}$ 。
\begin{subproblem}
求电流在通电导线外产生的总磁场的空间分布；
\end{subproblem}
\begin{subproblem}
电流在原点附近产生的总磁场的近似表达式，保留至线性项；
\end{subproblem}
\begin{subproblem}
将某原子放入磁阱中，该原子在磁阱中所受磁作用的束缚势能正比于其所在位置的总磁感应强度 $B_{tot}$ 的大小，即磁作用束缚势能 $V = \mu |B_{tot}|$ ， $\mu$ 为正的常量。求该原子在原点 $O$ 附近所受磁场的作用力；
\end{subproblem}
\begin{subproblem}
在磁阱中运动的原子最容易从 $x - y$ 平面上什么位置逸出？求刚好能够逸出磁阱的原子的动能。
\end{subproblem}
\end{problem}

